\section{Evaluasi}
\label{sec:evaluasi-bab12}

\subsection{Soal Pilihan Ganda}

\begin{enumerate}
    \item \lstasm|MOV AX, [BX]| memakai mode pengalamatan \dots
          \begin{enumerate}
              \item \textit{register}
              \item \textit{direct memory}
              \item \textit{indirect memory}
              \item \textit{immediate}
          \end{enumerate}

    \item \lstasm|MOV AX, [BX+SI+4]| memakai mode \dots
          \begin{enumerate}
              \item \textit{base relative}
              \item \textit{direct index}
              \item \textit{based indexed}
              \item \textit{indirect}
          \end{enumerate}

    \item Segmen default untuk \lstasm|MOV AX, [BP]| adalah \dots
          \begin{enumerate}
              \item \reg{CS} \item \reg{DS}
              \item \reg{SS} \item \reg{ES}
          \end{enumerate}

    \item Diketahui \reg{DS} = \hex{1000}, \reg{BX} = \hex{0200},
          \reg{SI} = \hex{0004}, dan \lstasm|MOV AX, [BX+SI+2]|.
          Alamat efektifnya adalah \dots
          \begin{enumerate}
              \item \hex{0204} \item \hex{0206}
              \item \hex{10206} \item \hex{10204}
          \end{enumerate}

    \item Kombinasi register yang legal untuk mode \textit{based
          indexed} adalah \dots
          \begin{enumerate}
              \item \code{[BX+BP]}
              \item \code{[SI+DI]}
              \item \code{[BX+DI]}
              \item \code{[AX+SI]}
          \end{enumerate}
\end{enumerate}

\subsection{Soal Uraian}

\begin{enumerate}
    \item Sebutkan tujuh mode pengalamatan 8086 dan berikan satu
          contoh untuk masing-masing.

    \item Jelaskan perbedaan mode \textit{direct memory},
          \textit{indirect memory}, dan \textit{base relative}.
          Berikan contoh penulisan untuk setiap mode.

    \item Sebutkan empat kombinasi register yang legal pada mode
          \textit{based indexed}. Hitung alamat efektif dan alamat
          fisik untuk \lstasm|MOV AX, [BX+SI+4]| bila
          \reg{DS} = \hex{2000}, \reg{BX} = \hex{0100}, dan
          \reg{SI} = \hex{0006}.

    \item Jelaskan segmen default untuk \reg{BX}, \reg{BP},
          \reg{SI}, dan \reg{DI}. Tentukan pula mode yang tepat untuk
          mengakses variabel tunggal, elemen array, medan record, dan
          elemen array of record.
\end{enumerate}

\subsection{Kisi-kisi Penilaian}

\begin{table}[H]
    \centering
    \caption{Kisi-kisi evaluasi Bab~\ref{chap:mode-pengalamatan}.}
    \label{tab:kisi-bab12}
    \small
    \begin{tabular}{@{}p{1.5cm}p{5.0cm}p{3.0cm}c@{}}
        \toprule
        \textbf{Sub-CPMK} & \textbf{Indikator yang diukur} &
        \textbf{Bentuk soal} & \textbf{Bobot} \\
        \midrule
        3.3 & Mengidentifikasi tujuh mode pengalamatan &
        Pilihan ganda 1, 2; Uraian 1, 2 & 35\% \\
        \addlinespace
        3.3 & Menerapkan kombinasi register dan menghitung alamat &
        Pilihan ganda 4; Uraian 3 & 30\% \\
        \addlinespace
        3.3 & Menentukan segmen default &
        Pilihan ganda 3; Uraian 4 & 20\% \\
        \addlinespace
        3.3 & Memilih mode untuk array dan record &
        Pilihan ganda 5; Uraian 4 & 15\% \\
        \bottomrule
    \end{tabular}
\end{table}

\begin{checklistbox}
\vspace{1.1em}
\textbf{Checklist Pencapaian Kompetensi
Bab~\ref{chap:mode-pengalamatan}}

Seorang mahasiswa dinyatakan kompeten pada Sub-CPMK 3.3 apabila:
\begin{itemize}
    \item[\checked] dapat menyebutkan tujuh mode pengalamatan;
    \item[\checked] dapat mengenali mode dari penulisan operand;
    \item[\checked] dapat membedakan \textit{direct} dan
          \textit{indirect};
    \item[\checked] dapat menggunakan \textit{base relative} dan
          \textit{direct index};
    \item[\checked] dapat menggunakan \textit{based indexed};
    \item[\checked] dapat menghitung alamat efektif dan alamat fisik;
    \item[\checked] dapat menentukan segmen default.
\end{itemize}
\end{checklistbox}
